%!TEX root = ../../main.tex

\chapter{Implementierung}
\label{ch:implementierung}

Das vorangegangene Kapitel~\ref{ch:konzeption} hat das Untersuchungsdesign entworfen und der Datensatz aus Kapitel~\ref{ch:datensatz} stellt die unveränderliche Datengrundlage bereit. Dieses Kapitel beschreibt, wie dieses Design in ein lauffähiges Softwaresystem überführt wurde, mit dem sich die vollständige Matrix aus Modell, Prompting-Strategie und \ac{ETL}-Aufgabe automatisiert durchführen und einheitlich bewerten lässt. Zunächst wird die modulare Systemarchitektur und die Versuchsumgebung dargestellt (Abschnitt~\ref{sec:architektur}), gefolgt von den eingesetzten Technologien und Bibliotheken (Abschnitt~\ref{sec:technologien}). Anschließend werden die Umsetzung der klassischen Vergleichspipeline (Abschnitt~\ref{sec:klassische-pipeline}), die Integration der \acp{LLM} über eine einheitliche Schnittstelle (Abschnitt~\ref{sec:llm-integration}) sowie die konkrete Realisierung der Prompting-Strategien (Abschnitt~\ref{sec:prompting-umsetzung}) beschrieben. Den Abschluss bildet das automatisierte Mess- und Auswertungsverfahren (Abschnitt~\ref{sec:messverfahren}), das die in Abschnitt~\ref{sec:evaluationskriterien} definierten Kriterien technisch operationalisiert und damit die Grundlage für die Durchführung in Kapitel~\ref{ch:durchfuehrung} schafft.

\section{Systemarchitektur und Versuchsumgebung}
\label{sec:architektur}

Das System ist als modulares Python-Projekt aufgebaut, dessen Komponenten entlang der einzelnen Verarbeitungsschritte getrennt sind und über klar umrissene Schnittstellen zusammenwirken. Diese Trennung folgt dem Ziel, die einzelnen Faktoren des Versuchsdesigns unabhängig voneinander austauschen zu können, ohne den übrigen Ablauf zu verändern, und erleichtert zugleich die Nachvollziehbarkeit der Ergebnisse \autocite{nagabhyruBridgingTraditionalETL2022,boulahiaMulticriteriaEvaluationResearch2024}. Ein Datenmodul kapselt die Erzeugung des synthetischen Datensatzes, die Referenzlösungen sowie die Definition der neun Aufgaben, ein Modellmodul stellt die Anbindung der \acp{LLM} bereit, ein Modul für die klassische Pipeline enthält die regelbasierte Vergleichslösung, ein Experimentmodul steuert den Ablauf und die Ausführung des erzeugten Codes, ein Evaluationsmodul bewertet die Ergebnisse, und ein Auswertungsmodul bereitet die erhobenen Messdaten für die Analyse auf.

Der grundlegende Datenfluss folgt für jede Konfiguration demselben Muster: Ausgehend vom synthetischen Datensatz und der Aufgabenbeschreibung wird ein Prompt erzeugt, an das jeweilige Modell übergeben und dessen Antwort weiterverarbeitet. Entsprechend der in Abschnitt~\ref{sec:forschungsdesign} getroffenen Gestaltungsentscheidung erzeugt das Modell dabei nicht das Ergebnis unmittelbar, sondern ein ausführbares Verarbeitungsskript, das anschließend deterministisch auf den Daten ausgeführt wird. Das so entstandene Ergebnis wird als tabellarische Ausgabedatei gespeichert, gegen die bekannte Soll-Lösung geprüft und gemeinsam mit allen erhobenen Kennzahlen in einem unveränderlichen Ergebnisdatensatz protokolliert \autocite{khanOverviewETLTechniques2024,zhaoSurveyLargeLanguage2026}. Sämtliche steuerbaren Randbedingungen sind in einer zentralen Konfiguration gebündelt, die unter anderem den Zufallskern für die Datenerzeugung, die konstant niedrig gehaltene Temperatur, die Obergrenze der Ausgabe-Token sowie die Anzahl der Wiederholungen je Konfiguration festlegt. Diese zentrale und explizit hinterlegte Parametrisierung ist eine wesentliche Voraussetzung für die in Abschnitt~\ref{sec:evaluationskriterien} als querschnittliche Dimension eingeführte Reproduzierbarkeit, da sie die Bedingungen jedes Durchlaufs eindeutig dokumentiert \autocite{zhaoSurveyLargeLanguage2026,hassaniRoleChatGPTData2023}. Die Ausführung erfolgt durchgängig auf demselben Rechner unter identischer Laufzeitumgebung. Die konkreten Angaben dazu finden sich in Abschnitt~\ref{sec:reproduzierbarkeit}.

\section{Eingesetzte Technologien und Bibliotheken}
\label{sec:technologien}

Als Implementierungssprache dient Python. Ausschlaggebend sind zwei nachprüfbare Gründe: Sämtliche untersuchten Modellanbieter stellen offizielle Python-Schnittstellen bereit, und die für die Datenverarbeitung benötigten Bibliotheken, allen voran Pandas, liegen dort in ausgereifter Form vor, wie auch die einschlägigen Übersichtsarbeiten zur automatisierten Datenverarbeitung durchgehend voraussetzen \autocite{mumuniAutomatedDataProcessing2025,coteDataCleaningMachine2024}. Die eigentliche Datenverarbeitung stützt sich sowohl in der klassischen Pipeline als auch im vom Modell erzeugten Code auf die Bibliothek Pandas, die sich als De-facto-Standard für tabellarische Verarbeitung etabliert hat. Ihre einheitliche Verwendung stellt sicher, dass die regelbasierte und die \ac{KI}-gestützte Lösung auf derselben technischen Grundlage arbeiten und damit unmittelbar vergleichbar bleiben \autocite{khanOverviewETLTechniques2024,nagabhyruBridgingTraditionalETL2022}. Sämtliche Zwischen- und Ergebnistabellen werden im spaltenorientierten Parquet-Format ausgetauscht, da dieses die Datentypen der Spalten mitführt und so einen typerhaltenden Abgleich zwischen erzeugter Ausgabe und Soll-Lösung ermöglicht, wie er für das Konsistenzkriterium aus Abschnitt~\ref{subsec:vollstaendigkeit} benötigt wird.

Die Anbindung der Modelle erfolgt über die jeweils offiziellen Softwarebibliotheken der Anbieter für die geschlossenen Modelle von OpenAI, Anthropic und Google sowie über eine lokale Laufzeitumgebung für das offen verfügbare Modell der Llama-Reihe. Diese heterogenen Schnittstellen werden, wie in Abschnitt~\ref{sec:llm-integration} beschrieben, hinter einer einheitlichen Abstraktion zusammengeführt, sodass der Versuchsablauf unabhängig vom konkreten Anbieter bleibt \autocite{zhaoSurveyLargeLanguage2026,2026AIIndex}. Ergänzend kommen etablierte Hilfsbibliotheken zum Einsatz: eine typsichere Konfigurationsverwaltung zur Bündelung der Versuchsparameter, eine Auszeichnungssprache zur versionierten Ablage der Prompt-Vorlagen sowie Werkzeuge zur Visualisierung und interaktiven Auswertung, mit denen die in Kapitel~\ref{ch:ergebnisse} dargestellten Tabellen und Abbildungen erzeugt werden. Als repräsentative Vertreter der jeweiligen Anbieter wird je ein aktuelles Modell eingesetzt, dessen exakte Versionskennung bei jedem Aufruf protokolliert wird, um der raschen Weiterentwicklung des Modellfeldes Rechnung zu tragen und die spätere Nachvollziehbarkeit zu sichern \autocite{2026AIIndex,zhaoSurveyLargeLanguage2026}.

\section{Umsetzung der klassischen ETL-Pipeline}
\label{sec:klassische-pipeline}

Die in Abschnitt~\ref{sec:baseline} als Vergleichsmaßstab definierte klassische Pipeline ist als regelbasierte Lösung umgesetzt, die dieselben neun Aufgaben ausschließlich mit fest vorgegebener Verarbeitungslogik und ohne Zugriff auf ein Sprachmodell löst. Für die klar formalisierbaren Aufgaben, etwa das Entfernen exakter Duplikate oder die Typkonvertierung, nutzt sie dieselben deterministischen Routinen, die auch der Erzeugung der Soll-Lösung zugrunde liegen. Für diese Aufgaben ist die explizit formulierte Regel zugleich die korrekte Lösung, sodass hier definitionsgemäß eine hohe Ergebnisqualität zu erwarten ist \autocite{walhaDataIntegrationTraditional2024,nagabhyruBridgingTraditionalETL2022}.

Von besonderer Bedeutung für die Aussagekraft des Vergleichs ist die bewusste Gestaltung der Pipeline bei den semantisch geprägten Aufgaben. Für die Vereinheitlichung der Länderangaben greift sie auf eine absichtlich begrenzte Zuordnungstabelle zurück, die zwar ausgeschriebene Ländernamen und Standardcodes abdeckt, jedoch keine Tippfehler, untypischen Abkürzungen oder umgangssprachlichen Bezeichnungen berücksichtigt. Nicht auflösbare Werte werden einheitlich als unbekannt gekennzeichnet. Für die Erkennung unscharfer Duplikate bildet sie einen normalisierten Schlüssel aus Name und E-Mail-Bestandteil und entfernt darüber Mehrfachnennungen, kann damit aber nur solche Abweichungen erfassen, die sich durch einfache Normalisierung angleichen lassen. Diese Heuristiken sind realistisch für ein rein regelbasiertes Vorgehen, stoßen jedoch genau dort an ihre Grenzen, wo ein inhaltliches Verständnis der Werte erforderlich wäre. Gerade an dieser Stelle wird der in Hypothese~H1 erwartete Unterschied zwischen regelbasierter und \ac{KI}-gestützter Verarbeitung messbar \autocite{hassaniRoleChatGPTData2023,coteDataCleaningMachine2024}. Da die Pipeline vollständig deterministisch arbeitet und identische Eingaben stets zu identischen Ausgaben führen, dient sie im weiteren Verlauf als stabiler und nachvollziehbarer Referenzpunkt, gegen den sich die Ergebnisse der Modelle einordnen lassen.

\section{Integration der LLMs in die Pipeline}
\label{sec:llm-integration}

Um die in Abschnitt~\ref{sec:modellauswahl} ausgewählten Modelle trotz ihrer unterschiedlichen Schnittstellen einheitlich ansprechen zu können, werden sie hinter einer gemeinsamen Abstraktion gekapselt. Jeder Anbieter implementiert dieselbe abstrakte Schnittstelle, die einen Aufruf entgegennimmt und eine standardisierte Antwort zurückliefert, in der neben dem erzeugten Text auch die exakte Modellkennung, der Verbrauch an Eingabe- und Ausgabe-Token sowie die Antwortzeit enthalten sind. Eine vorgeschaltete Erzeugungsinstanz liefert anhand eines Anbieternamens die passende Implementierung, sodass der Versuchsablauf vollständig unabhängig vom konkreten Modell bleibt und die Faktorachse Modell ohne Eingriffe in die übrige Logik variiert werden kann \autocite{zhaoSurveyLargeLanguage2026,2026AIIndex}. Diese standardisierte Antwort ist zugleich die Datenbasis für die spätere Effizienzbewertung nach Abschnitt~\ref{subsec:effizienz}, da sie die dafür benötigten Kennzahlen bei jedem Aufruf unmittelbar mitführt.

Der eigentliche Verarbeitungsschritt folgt dem Ansatz der Code-Generierung: Aus der Modellantwort wird der enthaltene Python-Code extrahiert und anschließend in einem eigenen Betriebssystemprozess mit fester Zeitbegrenzung ausgeführt. Ein Durchlauf gilt nur dann als erfolgreich, wenn der Prozess fehlerfrei endet und tatsächlich die erwartete Ausgabedatei erzeugt wurde. Andernfalls wird das Ergebnis als fehlgeschlagene Ausführung protokolliert. Da an dieser Stelle von einem Modell erzeugter Code ausgeführt wird, ist die Ausführung in einen separaten Prozess mit Zeitlimit ausgelagert. Dieses Vorgehen ist für die kontrollierte, lokale Durchführung der Experimente vertretbar, stellt jedoch bewusst keine vollständige Sicherheitsisolation dar und wäre für einen unbeaufsichtigten Produktivbetrieb entsprechend abzusichern \autocite{hassaniRoleChatGPTData2023,khanCognitiveETLUsing2025}.

Ein zentrales Element der Integration ist die Behandlung von Fehlern, die aufgrund der probabilistischen Natur der Modelle und der Nutzung externer Dienste unvermeidlich auftreten. Dabei wird bewusst zwischen zwei Arten von Fehlern unterschieden: Vorübergehende Störungen des Dienstes, die sich etwa in bestimmten Statuscodes oder Verbindungsabbrüchen äußern, werden als transient eingestuft und der betreffende Aufruf mit exponentiell wachsender Wartezeit mehrfach wiederholt. Ein inhaltlich fehlerhaftes, aber technisch wohlgeformtes Ergebnis, etwa ein Skript, das mit einem Programmfehler abbricht, wird hingegen nicht wiederholt, da es einen gültigen Messwert über die Fähigkeit des Modells darstellt und eine Wiederholung das Ergebnis verfälschen würde \autocite{zhaoSurveyLargeLanguage2026,hassaniRoleChatGPTData2023}. Ergänzend werden anbieterspezifische Besonderheiten aktueller Modelle abgefangen: Lehnt ein Modell einen Sampling-Parameter wie die Temperatur ab, wird dieser automatisch weggelassen und die Anpassung im Ergebnisdatensatz vermerkt, sodass solche Abweichungen bei der Interpretation der Reproduzierbarkeit sichtbar bleiben. Für jeden Aufruf werden schließlich sämtliche relevanten Angaben, von der Modellkennung über den Token-Verbrauch und die Kosten bis hin zum erzeugten Code und dem Ausführungsergebnis, vollständig protokolliert und bilden so die lückenlose Grundlage für die anschließende Auswertung.

Der in Abschnitt~\ref{sec:forschungsdesign} eingeführte zweite Ausführungsmodus, die direkte In-Context-Verarbeitung, ist über dieselbe einheitliche Modellschnittstelle umgesetzt, unterscheidet sich jedoch im Ablauf. Hier werden die Eingabedaten unmittelbar in den Prompt eingebettet und das Modell angewiesen, die transformierte Ergebnistabelle direkt in einem klar abgegrenzten Format zurückzugeben, das anschließend eingelesen und gegen die Soll-Lösung geprüft wird. Die Auswertung der Antwort ist dabei bewusst robust gegenüber am Ausgabelimit abgeschnittenen Antworten gestaltet, sodass auch unvollständige Ergebnisse noch verarbeitet und als solche erkannt werden. Da aktuelle Modelle standardmäßig einen internen Denkprozess ausführen, der das für die Ergebnistabelle verfügbare Ausgabebudget aufzehren würde, bevor die eigentlichen Daten erzeugt werden, wird dieses Verhalten für den Direkt-Modus deaktiviert und das Ausgabelimit angehoben \autocite{zhaoSurveyLargeLanguage2026,changEfficientPromptingMethods2024}. Anders als bei der wiederholten Code-Erzeugung wird jede Konfiguration in diesem Modus nur einmal ausgeführt, um den durch die eingebetteten Daten erhöhten Token-Verbrauch zu begrenzen. Eine Wiederholung erfolgt ausschließlich bei einem technischen Fehlschlag wie einem transienten Anbieterfehler oder einer leeren beziehungsweise nicht auswertbaren Antwort, nicht hingegen bei einem wohlgeformten, aber inhaltlich abweichenden Ergebnis, das einen gültigen Messwert darstellt.

\section{Umsetzung der Prompting-Strategien}
\label{sec:prompting-umsetzung}

Die in Abschnitt~\ref{sec:prompting-strategien-konzept} konzipierten Strategien sind als versionierte Vorlagen umgesetzt, die getrennt vom Programmcode in einem eigenen, menschenlesbaren Format abgelegt sind. Jede Vorlage besteht aus einer eindeutigen Kennung, einer Systemanweisung und einer Nutzervorlage mit klar benannten Platzhaltern, in die zur Laufzeit die Aufgabenbeschreibung sowie je nach Ausführungsmodus die Ein- und Ausgabepfade oder die eingebetteten Daten eingesetzt werden. Diese versionierte und vom Code entkoppelte Ablage stellt sicher, dass jede eingesetzte Prompt-Variante eindeutig dokumentiert, unverändert wiederverwendbar und damit nachvollziehbar bleibt, was eine wesentliche Voraussetzung für die Vergleichbarkeit und Reproduzierbarkeit der Ergebnisse ist \autocite{schulhoffPromptReportSystematic2025,changEfficientPromptingMethods2024}.

Alle drei Strategien folgen einem einheitlichen Grundaufbau und unterscheiden sich ausschließlich im Umfang der zusätzlich bereitgestellten Hilfestellung, sodass beobachtete Qualitätsunterschiede eindeutig auf die jeweilige Strategie und nicht auf abweichende Formulierungen zurückführbar sind. Die Zero-Shot-Vorlage beschränkt sich auf die reine Aufgabenanweisung und bildet damit die Referenzbasis. Die Few-Shot-Vorlage ergänzt ein vollständig ausformuliertes Ein- und Ausgabepaar. Dieses Beispiel bezieht sich bewusst auf eine andersartige Aufgabe auf einer separaten Beispieltabelle und nicht auf eine der neun evaluierten Aufgaben, damit es keine Musterlösung vorwegnimmt, sondern lediglich Stil und Format der erwarteten Lösung vermittelt \autocite{liuPretrainPromptPredict2023,schulhoffPromptReportSystematic2025}. Die Chain-of-Thought-Vorlage leitet das Modell an, seine Lösung in nachvollziehbaren Zwischenschritten zu entwickeln \autocite{sahooSystematicSurveyPrompt,schulhoffPromptReportSystematic2025}. Da alle Platzhalter für alle Strategien identisch belegt sind, bleibt der Vergleich über die Strategien hinweg fair und kontrolliert.

Für die Direktverarbeitung existiert zu jeder der drei Strategien eine eigene Vorlage, die dieselbe Hilfestellung bietet, aber statt eines Ausgabepfads die Eingabedaten selbst aufnimmt und statt eines Skripts die Ergebnistabelle anfordert. Diese Spiegelung ist Voraussetzung dafür, den Vergleich der beiden Ausführungsmodi in Kapitel~\ref{ch:ergebnisse} unter gleichen Prompt-Bedingungen zu führen, statt einen ausgewählten Prompt des einen Modus gegen einen festen Prompt des anderen zu stellen.

\section{Mess- und Auswertungsverfahren}
\label{sec:messverfahren}

Die in Abschnitt~\ref{sec:evaluationskriterien} definierten Kriterien werden durch eine automatisierte Bewertungsroutine operationalisiert, die die erzeugte Ausgabedatei mit der bekannten Soll-Lösung abgleicht. Um zu berücksichtigen, dass programmatisch erzeugte Ergebnisse korrekte Werte häufig in abweichender Zeilenreihenfolge liefern, werden beide Tabellen vor dem Abgleich einheitlich sortiert und damit reihenfolgeunabhängig verglichen. Die Genauigkeit wird als Anteil der zellweise übereinstimmenden Werte bestimmt, die Vollständigkeit aus dem Vorhandensein der erwarteten Spalten und der korrekten Datensatzanzahl abgeleitet und die Konsistenz über die Übereinstimmung der Datentypen mit dem Zielschema erfasst \autocite{ridzuanReviewDataQuality2024,boulahiaMulticriteriaEvaluationResearch2024}. Entsprechend der in Abschnitt~\ref{subsec:genauigkeit} getroffenen Unterscheidung sind beide Ausprägungen der Genauigkeit implementiert: Die strenge Variante bricht den Abgleich ab, sobald Spaltenmenge oder Datensatzanzahl abweichen. Die abgeglichene Variante, die der Auswertung zugrunde liegt, verlangt lediglich ein gültiges Zielschema, ordnet die Tabellen über ihre Schlüssel beziehungsweise die sortierte Reihenfolge einander zu und verbucht überzählige oder fehlende Datensätze als Abweichung. Numerische Spalten werden dabei mit einer Toleranz verglichen, sodass Rundungsunterschiede nicht als inhaltlicher Fehler zählen, während Textspalten exakt abgeglichen werden. Fehlt die strukturelle Grundlage vollständig, weil eine erwartete Spalte gar nicht vorhanden ist, wird das Ergebnis als Schema-Bruch gekennzeichnet und getrennt von einem Programmabbruch ausgewiesen. Die Effizienzkennzahlen aus Zeit, Token-Verbrauch und Kosten stammen unmittelbar aus der standardisierten Modellantwort und werden ergänzend zur inhaltlichen Bewertung erfasst \autocite{changEfficientPromptingMethods2024,2026AIIndex}.

Jeder einzelne Durchlauf wird als eigenständiger, unveränderlicher Ergebnisdatensatz abgelegt, der die vollständige Konfiguration, die erhobenen Qualitäts- und Effizienzkennzahlen, die exakte Modellkennung sowie den erzeugten Code umfasst. Diese lückenlose Protokollierung erlaubt es, jede spätere Auswertung ohne erneute Modellaufrufe aus den gespeicherten Datensätzen nachzuvollziehen, und bildet zugleich die Grundlage für die Erhebung der Reproduzierbarkeit, die sich aus der Streuung der Kennzahlen über die wiederholten Ausführungen einer Konfiguration ergibt \autocite{zhaoSurveyLargeLanguage2026,hassaniRoleChatGPTData2023}. Eine gesonderte Auswertungsschicht führt die einzelnen Datensätze zusammen und klassifiziert die Ergebnisse nach ihrem Zustand, um technische Fehlschläge von inhaltlichen Abweichungen zu trennen. Da jeder Lauf mit einem Zeitstempel im Dateinamen abgelegt wird und ein wiederholter Aufruf derselben Konfiguration, etwa nach einem Abbruch der Gesamtausführung, somit einen zweiten Datensatz erzeugt, statt den ersten zu überschreiben, wird je Konfiguration einschließlich ihrer Wiederholungsnummer der jeweils jüngste Datensatz herangezogen. Es handelt sich dabei ausdrücklich nicht um eine Auswahl nach Ergebnisqualität: Unterschieden wird allein nach dem Zeitpunkt der Ausführung, sodass eine unvollständig abgebrochene Gesamtausführung keine veralteten Werte in die Auswertung trägt. Auf dieser Grundlage werden die in Kapitel~\ref{ch:ergebnisse} dargestellten Vergleiche, Übersichten und detaillierten Abweichungsanalysen erstellt. Damit ist die technische Umsetzung des Untersuchungsdesigns abgeschlossen, und die tatsächliche Durchführung der Experimente sowie die dabei erhobenen Messdaten sind Gegenstand des folgenden Kapitels~\ref{ch:durchfuehrung}.
